\section{Discussion}
\label{sec:discussion}

\subsection{Why the effects occur}
\label{sec:why}

\emph{Stickiness is set by the cost distribution.} A policy that never moves
pays the cost of its cell less the cost of $\cstar(q)$, so its relative tax
is the ratio of the two costs less one. All three pools share a shape that
makes that ratio large: the top rung combines the largest model with its
costliest setting, while the mass of $\cstar(q)$ sits at the other end, among
the cheapest one or two options. The tax therefore follows how far the top
rung's price sits above the options that answer most prompts, which is 14
times locally, nine on the public pool and five on the hosted ladder. The
statement is exact, and follows from the decomposition \eqref{eq:decomp}.

\begin{proposition}[Stickiness of a fixed configuration]
\label{prop:sticky}
Let a policy use one cell $c$ for every prompt and never move, as P3 does with
$c = c_8$, and let $\mathcal{Q}$ be the servable prompts. Then, for every
$q \in \mathcal{Q}$, $\Esc(\pi, q) = 0$ and
$\Tax(\pi, q) = \Stk(\pi, q) = \cost(c, q) - \cost(\cstar(q), q)$, so that
\[
  \frac{\mathbb{E}_{\mathcal{Q}}\,[\Tax(\pi, q)]}
       {\mathbb{E}_{\mathcal{Q}}\,[\cost(\cstar(q), q)]}
  \;=\; \frac{\mathbb{E}_{\mathcal{Q}}\,[\cost(c, q)]}
             {\mathbb{E}_{\mathcal{Q}}\,[\cost(\cstar(q), q)]} \;-\; 1 .
\]
\end{proposition}
\begin{proof}
The policy makes one attempt, so $m = 1$ in \eqref{eq:decomp}: the sum in
$\Esc$ is empty and $a_m = c$. Taking means over $\mathcal{Q}$ and dividing by
the mean cost of $\cstar(q)$ gives the ratio.
\end{proof}
The relative tax of a policy that never moves is therefore the ratio of its
cell's mean cost to that of the minimum sufficient configuration, less one,
whatever the success rate of the cell. The relative taxes measured for P3,
12.96, 7.96 and 4.35 on the three pools, correspond to ratios of 13.96, 8.96
and 5.35, which are the figures above.

\emph{Escalation is small when the ladder's cost order matches the order of
sufficiency.} Each failed rung is paid in full, so a climb is cheap only if
its early rungs are both cheap and likely to succeed. The local P1 ladder has
that property: its first four rungs are the four cheapest cells, and they are
$\cstar(q)$ for 89.2\% of prompts. P2 breaks it at its second rung, where
thinking at 0.6B costs more than any larger non-thinking cell and answers
fewer prompts than any of them. The same property explains the other pools:
the hosted ladder's cheapest cell suffices for three queries in four, while
on the public pool 29.8\% of queries are answered most cheaply by a model the
frontier omits, so the climb cannot stop at $\cstar(q)$.

\emph{The costliest configuration is not the most accurate.} This holds in
every pool we measured on tasks with a checkable answer. Where we can see a
mechanism, it is the generation budget: between 16\% and 20\% of local
thinking generations, and 18.9\% of \texttt{gpt-5-nano}'s at high effort, run
out of tokens before they answer, and a truncated answer fails. A budget is
part of any deployed configuration, so sticky use of the top rung pays twice,
once in cost and again in the answers it loses. Open-ended requests are the
exception: there the cap rarely binds and $c_8$ is the most accurate cell,
yet sticky use of it still pays 5.39 times the minimum sufficient energy.

\emph{Thinking is priced by length, and length varies by family.} Energy per
output token barely moves with the mode, so the mode premium is the ratio of
output lengths, scaled by the change in accuracy. Plain mathematics answers
already write out their working, so thinking lengthens them about five times
and the energy ratio stays between 4.9 and 6.1, while plain code and
multiple-choice answers are shorter, thinking lengthens them 7 to 25 times,
and their ratios reach 15.8 and 15.2. The same accounting shows why a larger
model beats thinking below 8B: a size step keeps the output about as long and
raises energy per generation by 1.4 to 2.7 times, while a mode step keeps the
model and makes the output seven or eight times longer.

\emph{Routing too high is mostly a failure too.} Of the 968.6\,J premium for
routing one tier too high, only 167.5\,J is the higher tier's own excess
cost; the other 801.1\,J is recovery, because a step up in energy order is
not always a step up in accuracy and the step up fails on 24.1\% of prompts.
Against overshoot alone, routing too low would cost 15.7 times as much; the
factor of 2.71 in Section~\ref{sec:results-asymmetry} is lower because the
grid is not monotone, and it also depends on recovering at $c_8$.

\emph{A router loses to the frugal climb in recovery, not in prediction
alone.} P1's failures are cheap because its early rungs are the cheapest
cells. A router's failure is expensive because it recovers at $c_8$, which
costs eight times $c_7$ and answers fewer prompts. A router can avoid that
only by routing high, and the selected threshold does: the best router sends
three answerable prompts in four to a costlier cell than they need, so its
tax is almost all stickiness. Recovering by climbing from the routed cell
brings every router within 6\% of P1. The oracle's 32\% saving over P1
therefore needs a router that is rarely wrong, and the best router we trained
picks $\cstar(q)$ exactly for one prompt in seven. The dependence of the
threshold on the recovery rule can be stated exactly.

\begin{proposition}[Break-even threshold]
\label{prop:tau}
Let a router choose between a cheap cell $a$ and a safer cell $b$ with
$\cost(a, q) < \cost(b, q)$, and let $p_a$ and $p_b$ be their success
probabilities for $q$. Let $R_x > 0$ be the cost the recovery rule adds when
$x$ fails, and $\Delta = \cost(b, q) - \cost(a, q)$. Under oracle detection,
routing to $a$ has the lower expected cost if and only if
\[
  p_a \;\geq\; 1 - \frac{\Delta + (1 - p_b)\,R_b}{R_a}.
\]
When $b$ is certain to succeed, the break-even probability is
$\tau^{*}_a = 1 - \Delta / R_a$, which increases with $R_a$.
\end{proposition}
\begin{proof}
Routing to $x$ costs $\cost(x, q) + (1 - p_x) R_x$ in expectation. The cost of
$a$ is at most that of $b$ exactly when
$(1 - p_a) R_a \leq \Delta + (1 - p_b) R_b$; dividing by $R_a$ gives the
condition, and setting $p_b = 1$ gives $\tau^{*}_a$.
\end{proof}
A router should therefore demand more confidence before routing down the
dearer its recovery is, which is the sense in which the threshold belongs to
the recovery rule and not to the router. Given a failure, recovery at $c_8$
costs 3{,}432, 3{,}726 and 4{,}486\,J on average after a failure at $c_1$,
$c_3$ and $c_5$, over the 1{,}200 gradable prompts, while climbing the P1
ladder costs 1{,}367, 1{,}933 and 3{,}422\,J; after a failure at $c_7$ the two
coincide, at 4{,}834\,J. Climbing is the cheaper recovery wherever the two
differ, so the proposition predicts a lower threshold under climbing. The
selected thresholds agree: in all 46 runs, one per router and seed with the
cascade included, the energy threshold chosen on validation was lower under
climbing than under recovery at $c_8$, by a median of 0.55, and the same held
for latency; in tokens, where $c_8$ costs little more than the cells below it,
it fell in 28 runs, was unchanged in 15 and rose in 3. The proposition fixes
the direction and the dependence on $R_a$ but not the size of the shift, which
also reflects how well the router's probabilities are calibrated: taking the
safer cell to be certain, the break-even between adjacent non-thinking cells
is at least 0.93 under either rule, above every threshold selected for the
climbing rule.

\emph{Router cost does not separate the families on this grid; training data
does.} Every router's own energy is under 1\% of what it routes, so charging
it leaves the energy ranking almost untouched, and the order comes instead
from how well each family learns from 716 training prompts. A linear model on
$n$-grams or frozen embeddings needs little data, and the fine-tuned encoders
came close to it only when cost-sensitive weights focused training on the
errors that cost most. A 0.6B language model prompted to name the cheapest
sufficient configuration routes low, and on this grid routing low is the
expensive error. Only in tokens, where R3 pays for reading its own prompt,
does charging a router reorder the families. Generating first did not pay for
the cascade either: its checks cost 9.7\,J each, so verification is not what
makes it expensive; the information is weak, since a small model checking its
own answer accepts a wrong one on a third of answerable prompts.

\subsection{Limitations}
\label{sec:limitations}

\emph{The comparator is constructed.} The manual policies of
Table~\ref{tab:policies} are not observed behaviour. We have no logs of how
users of language-model services move between configurations, so the ladders,
the starting distribution of P4 and the detection probabilities are all
assumptions. We limit the damage by reporting a family of policies that
ranges from frugal to wasteful, by quoting headline figures under oracle
detection, in which the manual policies answer the most prompts, and by
sweeping the detection probabilities instead of fixing them. The tax we
report is therefore the tax these policies would pay, not a measurement of
what real users pay, and it is the part of the study that most needs evidence
from outside it. Two of these choices favour the frugal climb, the policy the
routers fail to beat: oracle detection assumes a user who recognises every
wrong answer, which is generous on mathematics and code and implausible on
open-ended work, and the climb makes the most attempts at which such a user
could be fooled, losing 12.3 points at the illustrative detection point.
Model loading is also uncharged, and the climb changes model more often than
any router.

\emph{Attempts are independent single-turn requests, and loading is not
charged.} Each rung is costed by \eqref{eq:policycost} as a fresh request
carrying only the prompt, and succeeds or fails as that cell did during
labelling. A user who escalates within one conversation also resends the
earlier answers, which lengthens every retry, and the stronger configuration
may be helped or misled by seeing a failed attempt; neither effect is
modelled, and the first would make escalation more expensive than we report.
On an 8\,GB GPU that holds one model at a time, a user who escalates across
sizes, or a router that sends consecutive queries to different sizes, also
pays the time and energy of loading a model, which our costs understate most
for whichever policy or router changes size most often.

\emph{Labels come from one sample under fixed caps.} With $k = 1$, each
$s(q, c)$ is a single draw of a stochastic outcome, and $\cstar(q)$ inherits
its noise: on the noise subset a cell's label differs between two samples
8.9\% of the time, and the single-sample $\cstar(q)$ matches the majority of
three for 84.0\% of prompts. The bootstrap resamples prompts with their
labels held fixed, so its intervals do not include that uncertainty. The
labels also depend on the caps. A thinking answer cut off at 4{,}096 tokens
counts as a failure, which is the right rule for deployment but makes
thinking cells look worse than they would with a longer budget. That budget
is below what Qwen3's own guidance suggests for thinking mode, and it
truncates 16 to 20\% of thinking generations, so it shapes three results that
would otherwise read as properties of the models: that $c_8$ answers fewer
prompts than $c_7$, that thinking is never the cheaper way to buy accuracy,
and part of P3's tax. The rerun bounds the effect rather than removing it: a
longer cap moved $\cstar(q)$ for 6.4\% of the capped generations rerun, and
27.9\% of them reached the 8{,}192-token cap too. What the grid shows is how
configurations behave under a deployment budget, not how the models behave
without one.

\emph{Grading is imperfect.} Mathematics answers are compared after
extraction and normalisation, which can reject a correct answer written in an
unusual form; an unparseable multiple-choice answer counts as a failure; and
code is judged by the datasets' own unit tests, which can pass an incorrect
function. The open-ended family depends on a judge that agrees with a human
only moderately, at $\kappa = 0.42$, and passes more answers than either the
human or \texttt{gpt-5-mini}; its results are reported separately and with
that caveat, and every other result covers the gradable families.

\emph{The energy boundary is the GPU board.} NVML reports GPU board power, so
our joules exclude the processor, system memory, storage and display, all of
which draw power while a query is served. Each query's energy is also
apportioned rather than measured, and CPU routers are charged an upper bound
rather than a measurement, which works against R1 in any comparison with a
router measured on the GPU.

\emph{Scope.} Every cell is a Q4\_K\_M quantisation of a Qwen3 model, so size
comparisons are between quantised models; the largest cell has 8B parameters;
thinking is either on or off, although Qwen3 also accepts a budget; and all
measurements come from one laptop and one GPU, so we do not claim the figures
transfer to other hardware. The gradable families come from public
benchmarks, in English and single-turn, which may appear in the models'
training data and are more uniform than live traffic---one reason the corpus
includes an open-ended family. The hosted arm is priced rather than metered,
from one provider's list prices on one date. The comparators are also
simplified: the cascade follows AutoMix without its meta-verifier and climbs
only the P1 ladder, and the prompted R3 routers read their choice from
next-token probabilities without training. Finally, the data are thin for the
comparison between router families---716 training prompts, 243 test prompts
and labels from one sample per cell---so their ranking should be read as what
these families achieve at this scale, not as a general ordering.

\subsection{Implications}
\label{sec:implications}

\emph{Defaults decide most of the tax.} Stickiness is between 74\% and all of
the tax for the policies that often start high, and in all three pools most
prompts are answered by the cheapest options. A service that opens on its
smallest adequate model, with extended reasoning off and a one-step-up
control in view, puts its users on the frugal climb, which spends 1.5 times
the oracle's energy per solved prompt against 5.7 times for leaving the
largest thinking cell selected. Turning reasoning on by default is the harder
choice to defend: within the cap it lost accuracy at 4B and 8B and was never
the cheaper way to buy accuracy on the gradable families, and at
\texttt{gpt-5} high effort answered 4.1\% of queries that minimal effort
missed, at 7.26 times the price, moving overall accuracy by a tenth of a
point. Reasoning is worth paying for on a minority of prompts, which argues
for choosing it per prompt rather than leaving it selected. The default moves
carbon as well as cost.

\emph{A user without a router should change size before mode.} The cheapest
habit on this evidence is to start small with thinking off, move up a size
when an answer fails, and turn thinking on only when the largest size without
it has failed. That advice depends on recognising a wrong answer: a user who
accepts one wrong answer in five keeps the climb's low cost but loses the
extra answers it buys. Prompts the user cannot check are therefore where a
one-shot router has most to offer, since its choice does not depend on their
judgment of the output.

\emph{Cost should be reported in the unit that is paid, per solved prompt.}
Output tokens price the mode axis to within about 15\% but miss the factor of
6.6 in energy per token between 0.6B and 8B, so in tokens P3 looks cheaper
than P2 and routing too low looks cheaper than routing too high, and energy
reverses both; a study that prices a mixed pool in tokens can therefore
favour exactly the misallocations that cost most. The wider pools also show
negative taxes for static policies that underpay by failing, so a tax, like a
cost, means little without the solved rate beside it.

\emph{Deploy a one-shot router against habits, not against a disciplined
climb.} Charged for its own energy, the best router spends 61\% less per
prompt than leaving the 8B thinking cell selected and 29\% less than users
who start at a uniformly random cell, answering more prompts than the first
and about as many as the second. It does not beat a user who climbs sizes
with thinking off and recognises every failure. Two conditions favour it:
when failures are not always recognised it answers 7.1 points more than the
climb, because it makes fewer attempts at which to accept a wrong answer, and
when its failures are recovered by climbing rather than by jumping to the
largest cell it comes within 6\% of the climb's energy. A service that routes
should therefore recover by climbing, and it gains most for users who cannot
check answers.
